% =================================================
% 性学系列 · 男性性健康
% 共享基座见 common.tex
% =================================================
\documentclass[12pt,UTF8,openany]{ctexbook}
% =================================================
% common.tex —— 性学系列 共享基座
% 导言区 / 配色 / TikZ 库 / ctexset / C 列型 / 字体 / 书眉页码
% 各册 % =================================================
% common.tex —— 性学系列 共享基座
% 导言区 / 配色 / TikZ 库 / ctexset / C 列型 / 字体 / 书眉页码
% 各册 % =================================================
% common.tex —— 性学系列 共享基座
% 导言区 / 配色 / TikZ 库 / ctexset / C 列型 / 字体 / 书眉页码
% 各册 \input{common.tex} 引用
% =================================================
\usepackage[top=2.5cm,bottom=2.5cm,left=2.6cm,right=2.6cm,headheight=2.0cm]{geometry}
\usepackage{amsmath}
\usepackage{amssymb}
\usepackage{array}
\usepackage{tabularx}
\usepackage{multirow}
\usepackage{booktabs}
\usepackage{enumitem}
\usepackage{xcolor}
\usepackage{graphicx}
\usepackage{tikz}
\usetikzlibrary{positioning,calc,arrows.meta,shapes.geometric,matrix}
\definecolor{myblue}{HTML}{185FA5}
\definecolor{myblueL}{HTML}{E6F1FB}
\definecolor{mygreen}{HTML}{3B6D11}
\definecolor{mygreenL}{HTML}{C0DD97}
\definecolor{myteal}{HTML}{1D9E75}
\definecolor{myamber}{HTML}{854F0B}
\definecolor{myamberL}{HTML}{FAEEDA}
\definecolor{myred}{HTML}{A32D2D}
\definecolor{myredL}{HTML}{FCEBEB}
\definecolor{mypurple}{HTML}{534AB7}
\definecolor{mypurpleL}{HTML}{EEEDFE}
\definecolor{mycoral}{HTML}{D85A30}
\definecolor{mycoralL}{HTML}{FAECE7}
\definecolor{mypink}{HTML}{D4537E}
\definecolor{mypinkL}{HTML}{F4C0D1}
\definecolor{mypinkM}{HTML}{ED93B1}
\definecolor{mygray}{HTML}{888780}

% 字体
\setCJKfamilyfont{hei}{SimHei}
\newcommand{\myhei}{\CJKfamily{hei}}

% 书眉与页码（印刷友好：中文页眉 + 居中页码）
\usepackage{fancyhdr}
\pagestyle{fancy}
\fancyhf{}
\fancyhead[L]{\small\myhei 性学系列}
\fancyhead[R]{\small\leftmark}
\fancyfoot[C]{\small\thepage}
\renewcommand{\headrulewidth}{0.4pt}
\fancypagestyle{plain}{\fancyhf{}\fancyfoot[C]{\small\thepage}\renewcommand{\headrulewidth}{0pt}}
\fancypagestyle{empty}{\fancyhf{}\renewcommand{\headrulewidth}{0pt}}

% 行距（零基础易读，略松）
\linespread{1.3}

% 超链接：印刷友好（黑色，电子版仍可点击）
\usepackage[colorlinks,linkcolor=black,citecolor=black,urlcolor=black]{hyperref}

\ctexset{
  chapter/name={第,章},
  chapter/number={\arabic{chapter}},
  section/format=\Large\bfseries,
}

\newcolumntype{C}{>{\centering\arraybackslash}X}
 引用
% =================================================
\usepackage[top=2.5cm,bottom=2.5cm,left=2.6cm,right=2.6cm,headheight=2.0cm]{geometry}
\usepackage{amsmath}
\usepackage{amssymb}
\usepackage{array}
\usepackage{tabularx}
\usepackage{multirow}
\usepackage{booktabs}
\usepackage{enumitem}
\usepackage{xcolor}
\usepackage{graphicx}
\usepackage{tikz}
\usetikzlibrary{positioning,calc,arrows.meta,shapes.geometric,matrix}
\definecolor{myblue}{HTML}{185FA5}
\definecolor{myblueL}{HTML}{E6F1FB}
\definecolor{mygreen}{HTML}{3B6D11}
\definecolor{mygreenL}{HTML}{C0DD97}
\definecolor{myteal}{HTML}{1D9E75}
\definecolor{myamber}{HTML}{854F0B}
\definecolor{myamberL}{HTML}{FAEEDA}
\definecolor{myred}{HTML}{A32D2D}
\definecolor{myredL}{HTML}{FCEBEB}
\definecolor{mypurple}{HTML}{534AB7}
\definecolor{mypurpleL}{HTML}{EEEDFE}
\definecolor{mycoral}{HTML}{D85A30}
\definecolor{mycoralL}{HTML}{FAECE7}
\definecolor{mypink}{HTML}{D4537E}
\definecolor{mypinkL}{HTML}{F4C0D1}
\definecolor{mypinkM}{HTML}{ED93B1}
\definecolor{mygray}{HTML}{888780}

% 字体
\setCJKfamilyfont{hei}{SimHei}
\newcommand{\myhei}{\CJKfamily{hei}}

% 书眉与页码（印刷友好：中文页眉 + 居中页码）
\usepackage{fancyhdr}
\pagestyle{fancy}
\fancyhf{}
\fancyhead[L]{\small\myhei 性学系列}
\fancyhead[R]{\small\leftmark}
\fancyfoot[C]{\small\thepage}
\renewcommand{\headrulewidth}{0.4pt}
\fancypagestyle{plain}{\fancyhf{}\fancyfoot[C]{\small\thepage}\renewcommand{\headrulewidth}{0pt}}
\fancypagestyle{empty}{\fancyhf{}\renewcommand{\headrulewidth}{0pt}}

% 行距（零基础易读，略松）
\linespread{1.3}

% 超链接：印刷友好（黑色，电子版仍可点击）
\usepackage[colorlinks,linkcolor=black,citecolor=black,urlcolor=black]{hyperref}

\ctexset{
  chapter/name={第,章},
  chapter/number={\arabic{chapter}},
  section/format=\Large\bfseries,
}

\newcolumntype{C}{>{\centering\arraybackslash}X}
 引用
% =================================================
\usepackage[top=2.5cm,bottom=2.5cm,left=2.6cm,right=2.6cm,headheight=2.0cm]{geometry}
\usepackage{amsmath}
\usepackage{amssymb}
\usepackage{array}
\usepackage{tabularx}
\usepackage{multirow}
\usepackage{booktabs}
\usepackage{enumitem}
\usepackage{xcolor}
\usepackage{graphicx}
\usepackage{tikz}
\usetikzlibrary{positioning,calc,arrows.meta,shapes.geometric,matrix}
\definecolor{myblue}{HTML}{185FA5}
\definecolor{myblueL}{HTML}{E6F1FB}
\definecolor{mygreen}{HTML}{3B6D11}
\definecolor{mygreenL}{HTML}{C0DD97}
\definecolor{myteal}{HTML}{1D9E75}
\definecolor{myamber}{HTML}{854F0B}
\definecolor{myamberL}{HTML}{FAEEDA}
\definecolor{myred}{HTML}{A32D2D}
\definecolor{myredL}{HTML}{FCEBEB}
\definecolor{mypurple}{HTML}{534AB7}
\definecolor{mypurpleL}{HTML}{EEEDFE}
\definecolor{mycoral}{HTML}{D85A30}
\definecolor{mycoralL}{HTML}{FAECE7}
\definecolor{mypink}{HTML}{D4537E}
\definecolor{mypinkL}{HTML}{F4C0D1}
\definecolor{mypinkM}{HTML}{ED93B1}
\definecolor{mygray}{HTML}{888780}

% 字体
\setCJKfamilyfont{hei}{SimHei}
\newcommand{\myhei}{\CJKfamily{hei}}

% 书眉与页码（印刷友好：中文页眉 + 居中页码）
\usepackage{fancyhdr}
\pagestyle{fancy}
\fancyhf{}
\fancyhead[L]{\small\myhei 性学系列}
\fancyhead[R]{\small\leftmark}
\fancyfoot[C]{\small\thepage}
\renewcommand{\headrulewidth}{0.4pt}
\fancypagestyle{plain}{\fancyhf{}\fancyfoot[C]{\small\thepage}\renewcommand{\headrulewidth}{0pt}}
\fancypagestyle{empty}{\fancyhf{}\renewcommand{\headrulewidth}{0pt}}

% 行距（零基础易读，略松）
\linespread{1.3}

% 超链接：印刷友好（黑色，电子版仍可点击）
\usepackage[colorlinks,linkcolor=black,citecolor=black,urlcolor=black]{hyperref}

\ctexset{
  chapter/name={第,章},
  chapter/number={\arabic{chapter}},
  section/format=\Large\bfseries,
}

\newcolumntype{C}{>{\centering\arraybackslash}X}


\begin{document}

% ===== 封面 =====
\begin{titlepage}
\centering
\vspace*{1.5cm}
\begin{tikzpicture}
\draw[myblue,line width=2pt] (0,0) circle (1.2);
\node[text=white] at (0,0) {\myhei\zihao{3}性};
\end{tikzpicture}
\vspace{0.6cm}
\noindent\rule{\textwidth}{0.4pt}
\vspace{0.5cm}
\myhei\zihao{5} 性学系列\,·\,家庭健康科普\,·\,全五册
\vspace{2.5cm}
\myhei\zihao{0} 男性性健康\par
\vspace{1.2cm}
\itshape\zihao{3} 男性身体的专项知识\par
\vspace{5.2cm}
\noindent\rule{\textwidth}{0.4pt}
\vspace{0.5cm}
\zihao{5} 性学系列\,·\,\today
\end{titlepage}

% ===== 免责声明 =====
\begin{center}\zihao{5}\itshape
本书为家庭健康科普读物，内容力求科学严谨，但\textbf{不替代专业医生的诊断与治疗}。\\
涉及个人健康问题，请以正规医疗机构的检查与建议为准。\\
全书以安全、自愿、相互尊重为底色，尊重多元与个体差异。\par\end{center}
\clearpage

% ===== 如何使用本书 =====
\chapter*{如何使用本书}
\addcontentsline{toc}{chapter}{如何使用本书}
\noindent\textbf{性学系列共五册：}
\begin{itemize}[leftmargin=2em]
\item \textbf{A 性健康基础} —— 认知、同意、防护、STI、避孕总论（建议先读）
\item \textbf{B 亲密实践与技巧} —— 前戏、体位、高潮引导、呼吸与盆底
\item \textbf{C 男性性健康} —— 男性解剖、ED/自慰、绝育、不育
\item \textbf{D 女性性健康} —— 女性解剖、妇科、不孕、人流、避孕、孕期/更年期
\item \textbf{E BDSM 与 Kink} —— 安全框架下的亲密表达
\end{itemize}
\vspace{0.6em}
\noindent\textbf{建议阅读路径（零基础）：} 先通读 A 建立认知与安全底线 → 按需读 C 或 D 了解自己的身体 → 再读 B 学技巧 → 兴趣到了再读 E。每章末有“一页要点”可快读。
\clearpage

% ===== 目录 =====
\tableofcontents
\clearpage
\chapter{男性视角的补充：同样重要的一半}
\label{chap:male}

前面各章更多是"如何保护自己"的口吻。但性是两个人的事，
男性有对等的权利、责任与困扰。这一章从男性角度补齐。

\section{同意是双向的：你也有权说不}

\begin{itemize}[leftmargin=2em]
  \item 前面讲"学会说不"，对男性同样成立。\textbf{对方有同意权，你也有}。
  \item 不必因为"是男人"就必须主动、必须想、必须持续。没准备好就说"我今天不想"，不必解释。
  \item 被对方施压（"你不配合就是不爱我"）和你施压一样，都是越界。
  \item \textbf{别把"不抗拒"当成"同意"}——明确、清醒、可撤回的肯定，才是同意。
\end{itemize}

\section{男性也会成为受害者}

这不是小事。WHO 数据显示，约\textbf{1/5 的女性、1/7 的男性}报告童年遭受过性侵。
男性遭遇性骚扰、性侵、亲密关系暴力的比例不低，但因为"男人就该扛着"的观念，
很多人选择沉默。

\begin{itemize}[leftmargin=2em]
  \item 遭遇侵害不是你的错，与"强不强壮"无关。
  \item 求助渠道对男性同样开放：拨打 110、12320，或联系医院、可信任的人。
  \item WHO 明确指出：僵化的"男性气概"规范会损害男性健康——鼓励冒险、酗酒、不愿求助。
        主动求助是勇敢，不是软弱。
\end{itemize}

\section{常见身体困扰：常见、可治疗、别乱补}

\subsection*{1. 勃起功能障碍（ED）是全身健康的信号}

\begin{itemize}[leftmargin=2em]
  \item 不同研究数字差异很大（如波士顿 MMAS 研究 40--70 岁约 52\%，
        科隆研究 30--80 岁约 19.2\%），共同结论是\textbf{随年龄上升}。
  \item 关键：ED 常与糖尿病、高血压、血脂异常等代谢/心血管问题相关。
        EAU 指南要求在诊治时评估血糖、血脂和睾酮——\textbf{它可能是身体给出的早期警报}，
        值得正规就医，而不是买补药。
  \item 治疗是一套成熟方案：先改善生活方式（运动、减重、控压控糖），
        一线药物为 PDE5 抑制剂，均需在医生指导下使用。
\end{itemize}

\subsection*{2. 早泄（PE）：先搞清楚是不是"病"}

\begin{itemize}[leftmargin=2em]
  \item ICD-11 已把它改称"早期射精"：指插入前或插入后极短时间内射精、
        几乎无控制感、持续数月并带来明显困扰。
  \item EAU 强调\textbf{不能只看时间}，还要看控制感、是否普遍、是否造成困扰。
        "正常变异"（时快时慢）和"主观性早泄"（时间其实正常但自认为快）都不算真正的疾病。
  \item 分\textbf{原发}（一直如此）与\textbf{继发}（后来才出现，可能与前列腺、激素、药物或心理因素有关）——继发的更要先找原因。
  \item \textbf{心态}：别把它当"个人缺陷"——它是\textbf{常见、可管理}的，伴侣的理解与沟通本身就是治疗的一部分。
  \item 真正需要时，行为训练、盆底肌锻炼、药物等都是有效手段，去泌尿外科/男科评估即可，\textbf{别独自硬扛或乱买补药}。
\end{itemize}

\textbf{治疗路径（从简到繁，多数轻度不需要药）：}

\begin{itemize}[leftmargin=2em,itemsep=0.3em]
  \item \textbf{行为法·停-起始法}：刺激到将射未射（约七八成兴奋）时停住，待射精冲动消退后再继续，如此重复数轮、最后一轮再完成。目的是让大脑重新识别"临界点"、建立"刹车"反射。自慰时即可练习，不必等到同房。
  \item \textbf{行为法·挤压法}：接近高潮时，用指腹轻压龟头下方（系带处）数秒，直到射精冲动消退。
  \item \textbf{盆底肌（PC 肌）锻炼}：即排尿时能"中途刹车"的那组肌肉。收缩保持约 5 秒、放松约 5 秒为一组，每组 10 次、每天 2--3 组，坚持数周见效；注意\textbf{单独收缩该肌、不憋气、不夹臀}。它同时改善勃起硬度与射精控制。
  \item \textbf{局部麻醉剂（延时喷剂/药膏）}：含利多卡因/丙胺卡因等，降低龟头敏感度，无系统性副作用；可能波及伴侣敏感度，注意用量、吸收后再行。
  \item \textbf{处方药（需医生评估）}：\textbf{达泊西汀}是目前唯一获批按需用于早泄的口服药，性生活前 1--3 小时服用、按需而非每日；另有部分 SSRI（如帕罗西汀/舍曲林）可超说明书用于早泄，需\textbf{每日服用、起效较慢（1--2 周}），有性欲下降、嗜睡等副作用。\textbf{切勿自行抓药，或与正在服用的抗抑郁药混用。}
\end{itemize}

\textbf{先排查可逆病因}：慢性前列腺炎、长期大量饮酒、焦虑与"表现压力"都可表现为早泄——\textbf{治原发病}往往比单纯"延时"更根本（本册《男性不育》章论及前列腺、《表现焦虑》一节论及压力）。

\subsection*{3. 别信"壮阳保健品"}

来路不明的补药、偏方成分不明，可能有害且耽误真正病因的诊治。有困扰就去看正规医院。

\section{男性在避孕与防病中的责任}

\begin{itemize}[leftmargin=2em]
  \item \textbf{安全套（避孕套）}：男性主导使用的工具，主动准备好、正确佩戴，是责任也是保护自己。
  \item \textbf{输精管结扎}：高效、永久性的男性避孕方式。CDC 数据：典型使用第一年失败率约
        0.15\%；术后 8--16 周需做精液检查确认无精子，确认前仍需用屏障避孕。
        它\textbf{不预防性病}。后悔率约 5\%，决定前要充分咨询。
  \item \textbf{HPV 疫苗男性也要打}：CDC 建议男女都在 11--12 岁常规接种、可 9 岁起，
        未接种者 26 岁前补种，27--45 岁与医生共同决策。可预防生殖器疣、肛门癌等。
  \item \textbf{乙肝疫苗}：有多个性伴或接受性病诊疗者应接种。
  \item 定期做性病筛查，同样是对自己和伴侣负责。
\end{itemize}

\section{"表现焦虑"多半是被教出来的}

\begin{itemize}[leftmargin=2em]
  \item 色情内容里的时长、强度、反应都是表演，拿它当标准必然焦虑。
  \item 一次性生活"不理想"非常常见——疲劳、压力、饮酒、情绪都会影响。
        偶尔失常不等于"不行了"。
  \item 越关注"表现"，越容易真的出问题（焦虑本身就是 ED 的常见原因）。
        把注意力放回连接、沟通和感受，反而更好。
  \item 真正需要警惕的是问题\textbf{持续出现}并带来困扰——那就去就医，而不是独自硬扛。
\end{itemize}

\section{给男性的一页清单}

\begin{enumerate}[leftmargin=2em,itemsep=0.3em]
  \item 我尊重对方的同意，也保有自己说"不"的权利。
  \item 我知道男性也可能成为受害者，遇到侵害会求助而不是硬扛。
  \item 我明白 ED/早泄常见且可治疗，ED 还可能是全身健康信号，会去正规医院而非买补药。
  \item 我会主动准备安全套，知道输精管结扎与 HPV 疫苗等选项。
  \item 我不拿色情内容当标准，不把一次性失常当成"不行"。
  \item 我不迷信"男人气概"，会表达情绪、会求助。
\end{enumerate}

\chapter{旧料新用：当伴侣遇到勃起困难，以及关于自慰的正常认知}
\label{chap:partner-ed}

\section{一、为什么选这一章}

本书前面讲过勃起功能障碍（ED），但都是从\textbf{当事人自己怎么办}的角度。本章换一个很少被正面讲的视角——\textbf{伴侣如何面对和支持}。这也是把旧书里那章《妻子如何面对丈夫阳痿与自慰》"旧料新用"的示范：它有很强的共情价值，但带着几处放到今天已过时的医学结论。下面先还原旧章说了什么，再用权威指南逐条修正，最后给出一版可直接并入新书的正文。

\section{二、旧章讲了什么（原样）}

\begin{itemize}
  \item \textbf{ED 定义}：阴茎无法勃起或无法维持足够硬度完成性交，持续较久；提到"约 10\% 成年男性"，随年龄增加。
  \item \textbf{成因}：分为\textbf{生理}（约 75\%，血管问题最常见——动脉供血不足/静脉回流过多，还有糖尿病、吸烟、药物副作用、神经/内分泌问题）和\textbf{心理}（约 25\%，紧张、表现焦虑、婚姻不和）。
  \item \textbf{主角是伴侣的心理冲击}：男方觉得"我不是个失败的男人吗？她会离开我吗？"；女方会误以为"是我没吸引力了/他有外遇了"——因为男方往往隐瞒、找借口回避亲密，反而让关系更冷。
  \item \textbf{治疗}：列了注射、真空泵、假体手术、心理辅导、激素、改变生活习惯等。
  \item \textbf{后半讲自慰}：引 1970 年调查（72\% 丈夫、68\% 妻子有自慰），说"婚姻中自慰完全正常"，是"快速的性松弛"，也是"对新奇的需求"。
\end{itemize}

\section{三、旧章的过时点（用今天权威指南修正）}

\begin{tabularx}{\textwidth}{lX}
\toprule
\textbf{旧章说法} & \textbf{今天怎么修正} \\
\midrule
用"\textbf{阳痿}"这个词 & 现代用"\textbf{勃起功能障碍（ED）}"。"阳痿"带强烈污名和羞辱意味，本身就会加重焦虑——名字换了，求助门槛就低了。 \\
治疗列表里\textbf{没有口服一线药}，把"口服治疗"写成只有育亨宾（Yohimbine）且"缺大型研究" & \textbf{EAU 指南明确：PDE5 抑制剂（西地那非/他达拉非等）是 ED 的一线治疗}（强推荐），起效可靠。育亨宾这类早已不是主流。这是旧章最大的信息缺口——很多人不知道"有药、且很有效"，白白忍了多年。 \\
把 ED 主要归为"心理/肾亏"式的个人问题 & \textbf{ED 常常是全身血管健康的早期信号}。EAU 指出 ED 与心血管疾病独立相关，可能是动脉粥样硬化的\textbf{首个临床表现}。出现 ED 该评估血压、血糖、血脂、吸烟等心血管风险——这不是"面子问题"，是\textbf{健康预警}。 \\
自慰被隐含为"婚姻不满足的补偿""怕被发现""精神不忠" & 现代共识：\textbf{自慰是正常、普遍、无害的性表达}，不需要羞耻，也不等于对伴侣不忠。只有当它变成强迫性、明显影响生活或关系时才需要关注。1970 年的数字虽旧，但"普遍存在"的结论成立。 \\
\bottomrule
\end{tabularx}

\vspace{0.5em}
\noindent\textbf{一句话：} 旧章的共情和"伴侣要一起面对"的内核对；过时的只是"病名污名 + 没提一线药 + 没上升到心血管预警 + 把自慰病理化"。把这些换掉，它就是一章很好的补充。

\section{四、重写后的正文（可直接并入新书）}

核心框架：\textbf{正名 $\rightarrow$ 成因与心血管预警 $\rightarrow$ 伴侣怎么做 $\rightarrow$ 现代治疗 $\rightarrow$ 自慰的正常化 $\rightarrow$ 何时就医}。

\subsection*{正名}
叫 ED，不叫"阳痿"；大多数男性一生中偶尔都会遇到，\textbf{持续 3 个月以上}才需评估。名字一改，求助门槛就低了。

\subsection*{成因与心血管预警}
血管问题最常见；糖尿病、高血压、吸烟、高血脂、部分降压/抗抑郁药都可能参与；心理因素常与生理因素\textbf{叠加}，不是"二选一"。\textbf{关键点：ED 常常是全身血管健康的早期信号}，出现 ED 该评估血压、血糖、血脂、吸烟——这不是"面子问题"，是健康预警。

\subsection*{伴侣四要三不要}
\begin{itemize}
  \item \textbf{要}：不指责、一起就医、把 ED 当健康问题（不是"爱不爱"的证明）、接受亲密不等于插入（拥抱/爱抚/相互自慰都算）。
  \item \textbf{不要}：不当场追问施压、不把它解读成"嫌弃我"、不替对方隐瞒回避。
\end{itemize}

\subsection*{现代治疗}
\textbf{PDE5 抑制剂是一线}（西地那非/他达拉非，需处方；\textbf{正在服硝酸酯类心血管药者禁用}）；生活方式改善（戒烟、运动、控糖控压）；无效再考虑注射/真空装置/手术；伴侣可一起参与咨询。旧章最大缺口就是"没提有特效一线药"，很多人因此白白忍了多年。

\subsection*{自慰的正常化}
正常、普遍、无害；是性欲差异的合理出口，不等于不忠；双方可坦诚沟通，让它成为关系的一部分而非秘密。旧章把自慰隐含成"婚姻不满足的补偿/精神不忠"，这是要纠正的。

\subsection*{红旗}
ED 可能是心血管病的首诊信号；\textbf{突发 ED 伴胸痛/气促要尽快就医}；持续困扰就找泌尿外科/男科。

\vspace{1em}
\noindent\textbf{一句话总结：} 保留旧章对人的关怀，换掉过时的医学结论——\textbf{ED 是健康信号而非面子问题、有一线药可用、伴侣要一起面对、自慰是正常的性表达}。

\chapter{男性生殖系统}

男性生殖器分为内生殖器和外生殖器。内生殖器包括生殖腺（睾丸）、输送管道（附睾、输精管、射精管和部分尿道）和附属腺体（精囊、前列腺和尿道球腺）。外生殖器包括阴囊和阴茎。

外生殖器

阴茎海绵体位于阴茎的背侧。左右并列构成阴茎的主体，尿道海绵体位于阴茎海绵体的腹侧，尿道贯穿其全长。其前端膨大为阴茎头，后端膨大称为尿道球。

海绵体为勃起组织，由许多海绵体小梁和腔隙组成，与血管相通，当海绵体腔隙充血时，阴茎变粗、变硬而勃起。阴茎海绵体和尿道海绵体外面共同由阴茎深、浅筋膜和皮肤包被。

阴茎的皮肤薄而柔软，富有伸缩性，皮肤至阴茎颈部游离向前延伸，形成包绕阴茎头的双层环形皱襞，称阴茎包皮，包皮的前端围成环形的小口称包皮口。在阴茎头尿道外口与包皮相连的皮肤皱襞称包皮系带。

内生殖器

1．睾丸：睾丸位于阴囊内，左、右各一。睾丸是稍扁的卵圆形器官，表面光滑。睾丸的表面覆盖一层浆膜，浆膜的深面是白膜，在睾丸后缘，白膜增厚形成睾丸纵隔，纵隔的结缔组织发出许多小隔，将睾丸分隔成许多个睾丸小叶，每个小叶内有1～4条高度盘曲的小管，称为精曲小管。精曲小管在小叶顶端汇集成精直小管，精直小管进入纵隔为睾丸网。精曲小管间的疏松结缔组织为睾丸间质。

精曲小管管壁由复层生精上皮构成。生精上皮由生精细胞和支持细胞组成。上皮下方有基膜，基膜的外侧有胶原纤维和肌样细胞，肌样细胞的收缩有利于精子的排出。

生精细胞包括精原细胞、初级精母细胞、次级精母细胞和精子。

精原细胞：在青春期前，精曲小管中仅有精原细胞。青春期开始，在脑垂体产生的促性腺激素作用下，精原细胞开始增殖发育。一部分继续保持原始状态，称作A型精原细胞，另一部分则发育转化为B型精原细胞，后者经过分裂分化为初级精母细胞。

初级精母细胞：初级精母细胞经过染色体复制、同源染色体配对、联会、等位基因交换等，完成第一次成熟分裂，形成两个次级精母细胞。

次级精母细胞：次级精母细胞迅速完成第二次成熟分裂，形成精子细胞。精子细胞的核型为23X和23Y。

睾丸产生精子，精子形似蝌蚪，分头、尾两部分。精子头部呈扁梨形。由高度浓缩的核和顶体组成，顶体内含多种水解酶。精子尾部是精子的运动器官。可分为颈段、中段、主段及末段四部分。

人类精原细胞发育成精子。大约需要64～75天。

2．附睾：完成精子形成的过程以后，精子首先进入输出小管并进入弯曲的睾道管；附睾的尾部是精子主要贮藏的地方。附睾壁是一层分泌上皮，具有不运动的纤毛柱状细胞。这些细胞分泌的液体除了营养功能以外精子在此成熟。精子在附睾中存活的时间要比在男性生殖道的其他部位长。在附睾中精子是不活动的，只是当它们和氧或一种能被代谢成乳酸的物质接触时才开始活动。

与整个精液相比，附睾分泌物的特点是钾的含量较高。一些精子的细胞内成分很可能是在附睾中获得的。附睾中吞噬细胞的细胞解体作用以及未射出的精子的降解在这个区域内进行，并可能为液体提供了某些成分。

3．输精管：附睾的尾部不太卷曲并发出一条很长的弯曲的导管，这就是输精管。输精管沿着睾丸的后缘向上，成为精索的一部分，最后横过膀胱进入前列腺的尿道。在进入前列腺之前，一条精液的主要分泌管加入输精管，两者形成了射精管，横过前列腺。在它进入前列腺之前管子变大，弯曲，称作为末端壶腹。在壶腹内是一层皱折的上皮细胞膜，成为精子的第二个贮存处，是一个在空间上和功能上复杂的系统。在输精管的近端有不活动的纤毛，但是在远端包括壶腹部却没有。壶腹是贮存和积累润滑精子的导管分泌物处。壶腹的分泌物是黄色的。

4．前列腺：人的前列腺并列于膀胱的颈部，直径大约4厘米，围绕着尿道和两条射精管，主要是由两个大的侧叶和较小的前叶、后叶和中叶所组成。二条射精管位于前列腺的中叶和侧叶之间。它们汇合时变小，并终止于前列腺尿道。射精管中排列着柱状的分泌上皮细胞，具有很多腺状的突起。这些突起伸入到前列腺内。前列腺的分泌物是碱性的，主要的成分是酸性磷酸酶、纤维蛋白溶解酶和钙。在射精的过程中，精子经前列腺过射精管，到尿道的前列腺部，一直到膜状尿道和阴茎尿道。在人类中无色的前列腺分泌物是偏酸性的（pH6.5）。具有一些很强的蛋白质分解酶。还含有大量的纤维蛋白分解酶；因此可以用1毫升的前列腺液在大约18小时内使50毫升已凝固了的血液液化。

5．尿道球腺和尿道腺：在人体中这些腺体很小，并由很多彼此紧密地联合在一起的小叶所组成。尿道球腺位于尿道的后部，分泌一种粘蛋白。这些物质是在精子之前被射出。在尿道的整个长度上有广泛分布的小腺体。

男性生殖生理

1．睾丸的功能：睾丸是男性的生殖腺，具有产生精子和分泌雄激素的功能。精子由睾丸精曲小管的精原细胞分阶段发育形成。雄激素由睾丸间质细胞合成和分泌。睾丸分泌的雄激素主要是睾酮，日分泌量约4～9毫克，50岁以上睾酮分泌量逐渐减少。

睾酮可以促进男性性器官的发育和副性征的出现，与卵泡刺激素共同刺激睾丸生精作用并维持椅子的成熟与活力；维持男性正常的性欲和性行为；促进肌肉和生殖器官的蛋白质合成，促进骨骼生长和钙磷沉积；促进红细胞生成。

睾丸功能由下丘脑-垂体-睾丸轴调节：下丘脑合成的促性腺激素释放激素经垂体门脉系统到达腺垂体。促进腺垂体促性腺激素细胞合成和分泌卵泡刺激素和黄体生成素。

睾丸分泌的睾酮在血中达到一定浓度后，可抑制下丘脑和腺垂体的分泌活动。致使促性腺激素释放激素、卵泡刺激素和黄体生成索的分泌减少。通过这种负反馈调节维持血中睾酮水平的相对稳定。

2．附睾的功能：

（1）附睾上皮有极强的吸收能力。附睾可分泌多种物质，与精子的营养和成熟有关。

（2）精子在附睾内停留15天左右，获得向前运动能力和受精能力。

（3）精子在附睾内经多种因素作用，向输精管方向移动。

（4）附睾尾部可暂时贮存少量精子。

3．附属腺体的功能：

（1）精囊分泌液呈碱性，可中和阴道酸性保护精子活动能力。

（2）前列腺分泌液呈弱碱性，可降低阴道酸度，有利于精子的生存和活动，影响精液的凝固和液化过程。

（3）尿道球腺分泌液有粘性，富含唾液酸蛋白，其主要作用是润滑尿道。

4．勃起与射精：

勃起与射精是一系列复杂的神经反射过程。勃起反射的基本中枢在脊髓骶段，此反射可由许多刺激引起，勃起明显受意识控制。阴茎的勃起主要是由于阴茎小动脉扩张，血窦充血膨胀所致，静脉血液也有助于勃起。射精首先由附睾，输精管平滑肌按一定顺序收缩。将精子排入后尿道，然后由坐骨海绵体肌和球海绵体肌收缩，将精液由尿道射出体外。

5．精液：

精液由精子和精浆两大部分组成。精子是男性生殖细胞。精浆主要由精囊和前列腺分泌液组成，约占精浆的90％以上，其余为附睾和尿道球腺分泌液。精浆有运送精子，中和阴道酸性，提供能源物质，参与精液的凝固和液化的作用。

总之，男性生殖功能包括：精子发生、精子成熟、精子排出以及通过性交把精液射入女性生殖道。

（1）精子发生。精子发生过程是在睾丸曲细精管内进行的。精原细胞经过有丝分裂、增殖、分化，再经过一次减数分裂以及变态等一系列复杂过程，最后生成成形精子。精子发生是雄性配子的产生过程，是在神经内分泌功能的严密调控下进行的。

（2）精子成熟。就功能和超微结构而言，在曲细精管中生成的精子尚未达到完全成熟，尚不具备受精能力。当精子进入附睾后，须在附睾中停留5～25天才能逐步成熟。此阶段精子无论在形态、大小、膜通透性、超微结构以及在代谢、耐寒和热性等方面均发生了不同程度的改变。在附睾中逐渐成熟的精子，其运动方式从原地摆动进而为转圈状运动，最后为螺旋式的向前运动，并获得使卵子受精的能力。精子在附睾中的以上变化称为精子的功能成熟过程。

（3）精子排出。精子与附属性腺分泌物共同组成精液，通过射精随精液一起排出体外，进入女性生殖道。

精子在穿入卵子并使卵子受精之前，必须在女性生殖道中度过一段时间，在这段时间，精子为使自己获得受精能力而发生的生理改变。此过程是精子获能的过程，无此过程，精子不能穿入卵子。

排入女性生殖道的精子，只有一个精子与卵子结合。精子与卵子结合后，在卵子外围便立刻形成一层膜，把其他精子拒之膜外。进入卵子的精子细胞核与卵细胞核接近，融合成合子，形成受精卵。





阴茎如何勃起


阴茎勃起，可能是由于思想或生理受到刺激所引致。也有可能是结合两种因素刺激所致。这些刺激包括视觉、味觉、听觉、幻觉或触觉，举例来说，男性只要想及性事或想到某种引起兴奋的物件或人物，便有可能勃起。脑部接收到刺激，便把信息传送到脊髓，然后再传送到阴茎去。

生殖器官部位受到刺激，亦会引致勃起。阴茎的感觉神经受到刺激后，把讯息转送到脊髓，再继续传送到大脑，然后传回阴茎去。一般情况是：在心理兴奋和阴茎同时受到生理上的刺激之下，向大脑发出讯息，阴茎周围的平滑肌开始放松，让血液供应增多，引致阴茎变大。

在正常勃起的情况下，需要从动脉。输送比日常高出六倍的血液到阴茎去。一般健康及性能力正常的年轻人，平均血液流量达每分钟九十毫升。在开始勃起之际，这个流量更增至每分钟二百七十毫升。





衡量勃起次数


现有一系列技巧可用以量度阴茎勃起的强度。其中一种早期的“用具”是用邮票在阴茎上缠上一个圆圈，假如翌晨发觉邮票的钉孔裂开，便可证明曾在睡梦中的勃起。这个简单的用具使用多时，直至发明Snap-Gauge量规。

Snap-Gauge量规用三张颜色胶纸制成，每张胶纸经过特别的设计，在随着某种程度的力量后便会爆裂。衡量的方法是在晚上把这种量规套上阴茎，翌晨检查有多少张胶纸爆裂，便可知究竟。

其后进一步发明了许多更加准确的量度装置。RigiScan便是其中一例。这种装置使用电池驱动，以细微的刻度记录睡中勃起的次数和力度。





我国男性一般睾丸大小


我国成年男子睾丸大小，范围在15到25号之间，15号相当于4.29厘米（纵径）×2.58厘米（横径）大小，25号相当5.01厘米（纵径）×3.01厘米（横径）大小，如果睾丸大小在此范围之内，应视作正常。睾丸长得太小，尤其离开上述正常标准距离很远者，性生活也许会妨碍。表现为阴茎勃起不够坚硬，射精没有力量等，但是多半人还是能够完成性生活，问题在于生育力，睾丸太小，其里面负责生产精子的曲细精管会发育不良，无力负担起生产精子的任务，往往导致不育。

\chapter{男性绝育手术}


男性绝育手术的概念


男性绝育术是将输精管结扎、切断或堵塞，使睾丸中产生的精子不能通过输精管排出，精液中无精子存在，因而达到避孕目的。

我国使用最早、最广泛的男性绝育手术为输精管结扎术。从1957年以来，随着我国计划生育工作的广泛开展，广大医务工作者对男性绝育术不断进行改革，先后创造发明了输精管银夹绝育术、输精管穿刺电凝绝育术、附睾注入药物绝育术、输精管注射粘堵术和输精管可复性栓堵绝育术等。后两种绝育手术目前处于世界领先地位，并在国内外广泛应用。





手术方法


传统的输精管结扎术常规作阴囊两侧切口，提出并游离输精管，切除1厘米，将两残端结扎。这样，由睾丸产生的精子便不能进入远端输精管道。由于精液中不再含有精子，因而不再能使女方怀孕。需要时，可以再作输精管吻合手术以恢复生育力。

直视钳穿法输精管结扎术手术简便、器械简单、手术时间短、出血少并易于掌握。直视钳穿法采用一种特制的输精管皮外固定钳将输精管连同被绷紧的阴囊皮肤一起套住，固定在阴囊皮下，用输精管分离钳在输精管最突出处穿入阴囊皮肤直至输精管组织并一次分开，将光裸的输精管提出，压挫并结扎，然后采用一种叫作筋膜包埋的技术处理结扎后的输精管残端。由于这种手术的一切操作均在直视下进行，避免了其他结扎方法要用锐器进入阴囊内或贯穿。

术中未用杀精药液作精囊灌注，术后一段时期内又未采用其他避孕措施，是输精管结扎术后近期再孕的一个常见原因。

输精管结扎后，精道远端和精囊内仍残留有活精子。以术后残留精子活力和活精子完全消失为标准，术后应避孕6个月或14～20次排精。如果条件许可。可作精液检查，证实确无精子，方可放弃避孕。在此之前应坚持避孕。

为了加速残余精子的消失，缩短术后避孕时限，我国著名泌尿外科专家吴阶平在1958年首次提出在术中向输精管远睾端低压缓慢灌注杀精药物，借助机械冲洗和药物的作用，杀死残余精子。在我国精囊灌注技术已作为输精管结扎术的一个必要步骤。被列入节育手术常规。国内外对精囊灌注的药物作了大量的研究，从临床应用结果看来，0.01％醋酸苯汞杀精效果最好；1/3000新洁尔灭和15％维生素C的效果也不错，而且取材方便，不失为好的杀精药物。

输精管结扎后，睾丸曲细精管仍然保持良好的生精状态，但与结扎前相比速度减慢。精子细胞仍然经历着与术前一样的生命周期；死亡后被机体吸收。在结扎以前，经附睾输精管排出的精子也只不过是其中的一小部分，绝大多数精子都在附睾内被重吸收。精子在附睾内的分解吸收始终维持着动态平衡。





手术后的并发症


输精管结扎术是一种安全、简便的门诊外科小手术，较严重的并发症十分罕见。手术后短期内大约半数男子发生阴囊局部瘀血、肿胀和疼痛，症状一般不严重，多在1～2周内自然消退。





输精管绝育术后的复通


男子行输精管绝育术后，需要恢复生育力时，可施行输精管吻合术，俗称再通术。由于结扎术或粘堵术仅是剪断、切除或损伤了极少的一段输精管，剩余部分再接通的可能性很大。至于栓堵术，只要将栓子取出，即可达到输精管复通。如果睾丸、附睾的功能正常，远端输精管仍通畅，则多可恢复生育能力。

目前，普遍采用显微外科技术进行操作，可减少组织损伤，有助于修复的成功。通常复通率可以达到95％，但经长期随访，妊娠率仅为75％。

术后精子出现的时候，早则1个月，晚者可在半年至1年，因此早期无精子还不能认为是手术失败。

当精子出现而妻子未能受孕者，可能有多方面的原因，应进行不孕症的详细检查。有报道说，在输精管结扎术后有50％以上的病例的血清中出现抗精子抗体，它是干扰受孕的因素之一。

阴囊壁等盲目性操作，对组织损伤极小，因此安全性高。手术伤口不用缝线，愈合迅速，术后血肿和感染的发生率较低。





手术的副作用及其处理方法


绝大多数人施行输精管结扎后无不良反应，但也有少数人会出现阴囊血肿、感染、痛性结节和附睾郁积症。

1．出血：适应证的选择失当或手术操作失当，都可以引起术后局部出血。出血一般发生在术后24小时内。根据出血部位不同，可分为阴囊皮下瘀血、精索血肿及阴囊血肿，需用促凝血药物，以抗生素预防感染，可用理疗方法促进积血的吸收。

2．感染：这是指因手术引起皮肤切口、精索或生殖道（附睾、输精管、前列腺及精囊）等部位的炎症，是手术后早期并发症之一。根据感染的部位不同，感染可分为阴囊切口感染、精索炎、前列腺炎和精囊炎几种类型。急性炎症应及时应用广谱抗生素并局部热敷、理疗。慢性炎症局部应用抗生素并用理疗，化脓者应及时切开引流。

3．痛性结节：术后3个月以上，自觉手术部位疼痛，检查输精管结节有压痛感，为痛性结节。这是输精管残端周围组织炎症引起的。感染型的应全身应用抗生素，局部封闭治疗，同时输以理疗；如长期不愈，应手术切除痛性结节。非感染型的，可局部封闭治疗配合理疗。输精管残端与阴囊壁粘连或有瘘管，应及时手术治疗。

4．附睾郁积症：术后3个月以上，自觉阴囊胀痛，检查附睾有肿大及压痛者，为附睾郁积症。这是局部血液循环不良、附睾重吸收功能不佳及局部炎症等原因引起。单纯性的，可用阴囊托局部理疗。长期不愈者应手术吻合输精管。如伴附睾炎，应加用抗生素治疗。

5．性功能障碍及心身疾病：受术者对输精管绝育术与性功能之间的关系缺乏科学的认识，或由于手术引起了出血、感染等并发症，都可能造成性功能障碍。常见的性功能障碍有性欲减退、阳痿、早泄和不射精等。有全身性疾病的应及时治疗，非器质性病变者应以精神治疗为主。





男性术前准备须知


1．手术前应由专业人员向受术者及其家属详细介绍输精管绝育术的原理，以取得密切配合，对预防术后产生的心理障碍有极大帮助。

2．受术者在术前1日应淋浴，并更换清洁内衣。

3．术前刮除阴毛，并用肥皂水、清水洗净阴囊部。





手术术后须知


1．术后当天静息，48小时内避免体力劳动和剧烈运动。

2．使用阴囊托48小时，如果觉得使用比不使用更舒服的话，还可以延长时间。

3．保持伤口部位清洁干燥48小时。避免撕扯、揉搓或刺激伤口。术后24小时不要洗澡，如果非洗不可，必须保持伤口干燥。

4．一周后可恢复性生活，如果没有什么不适感觉还可提早，若感觉不好则不应勉强。

5．在术后一段时间内仍有使妻子怀孕的可能。应采用其他形式避孕，如安全套等。最保险的途径是请医生在显微镜下检查精液，证实精液中确实已无精子，才能不用安全套。有条件时，在20次射精后应将精液标本送检。

6．如发生疼痛或不适，服用一般性的止痛剂。

7．术后阴囊区疼痛、肿胀、局部青紫均属正常现象，不必担忧。如果发生下述情况则应立即就诊：

发热（38℃以上）。

伤口渗血或有脓性分泌物。

严重的疼痛或肿胀。





其他一些男性节育方法


1．输精管注射粘堵术：该法采用阴囊皮肤穿刺法，将化学粘堵剂注入输精管内。粘堵剂为液体，注入输精管后能快速凝固，以阻塞输精管，阻止精子通过，从而达到绝育目的。由于这种绝育手术既不开刀，又不切断、结扎输精管，只要打一针就能达到绝育目的，消除了人们怕开刀的心理障碍，所以很受群众欢迎。

2．输精管可复性注射栓堵术：这是一种新型的绝育手术，也是用阴囊皮肤穿刺法将栓堵剂注入输精管内。栓堵剂在常温下为液体，注入输精管后很快凝固，形成一个有弹性的栓子，严密堵塞输精管，从而阻止精子通过，达到绝育目的。这种绝育方法的最大特点是可复性，如果需要再生育时，可以将栓取出即可恢复生育能力。

3．输精管银夹绝育术：这种方法是通过手术分离出一段输精管，然后用一种特殊的银夹（V形或U形）将输精管夹住、钳紧，造成管腔闭塞，以阻止精子通过。

4．输精管穿制电凝绝育术：这种手术是将特制绝缘针插入输精管腔，利用高频电流造成一小段输精管粘膜烫伤，以后逐渐结痂愈合。结痂后形成的瘢痕阻塞了输精管便可达到不育。


% =================================================
% C 册补充专题：男性作为性侵受害者（呼应 A 册《性侵犯与强奸》）
% 现代循证撰写
% =================================================

\chapter{男性作为性侵受害者（专科补充）}

\section{先把"男人也该扛住"放一边}
性侵不只是女性会遇到的事。\textbf{男性、男性之间的、以及异性恋男性被女性施加的}性侵都存在，
只是长期被"男人要主动、要硬、要能行"的刻板印象遮住，导致男性受害者\textbf{更少求助、更晚被发现}。
对男性来说，最大的障碍往往不是疼痛，而是\textbf{羞耻与被质疑}（"一个大男人怎么会""你当时怎么不反抗"）。
请把这两句先收下：\textbf{你是受害者，与你的性别、体力、是否"想发生"都无关；求助不是软弱。}

\section{男性特有的医学要点}
\begin{itemize}
  \item \textbf{生殖道 / 泌尿道损伤}：尿道或阴茎撕裂、血尿、射精痛，可能伤及尿道海绵体；出现"排尿困难、血尿、阴囊肿痛"要尽快到\textbf{急诊 / 泌尿外科}。
  \item \textbf{STI 暴露}：与 A 册一致——HIV（72 小时内 PEP 最有效）、梅毒、淋病、HPV、肝 B 等，需检测与随访。
  \item \textbf{意外致他人妊娠}：若为阴道 / 肛交且无保护，\textbf{伴侣}可能怀孕——这是男性受害者常被忽略、却又必须面对的一环。
  \item \textbf{会阴疼痛、盆底肌紧张}：部分男性事后出现会阴痛、勃起或射精不适，多与盆底肌痉挛 + 心理因素相关，可治疗。
\end{itemize}

\section{就医时怎么说（给"开不了口"的男性）}
\begin{itemize}
  \item 直接说：\textbf{"我可能遭遇了性侵，想检查一下伤情和感染风险，请帮我做相关检查。"}——一句话就够，不必展开细节。
  \item 可要求：体格检查 + 分泌物/尿液取样 + 血检（HIV/梅毒/乙肝）+ 必要时 STI 随访 + 伴侣妊娠评估的指引。
  \item 若不愿走法律程序，可只体检取证、保留证据，之后再决定是否报案；身体安全永远第一。
  \item 找\textbf{男性医生}或\textbf{同性别医护人员}陪同检查，能明显降低不适感。
\end{itemize}

\section{心理与社会支持}
\begin{itemize}
  \item 男性受害者同样可能出现 PTSD、闪回、对亲密关系退缩、性欲骤变（可能升高也可能降低，都正常）。
  \item 支持渠道：心理科 / 创伤聚焦咨询（TF-CBT、EMDR）、男性权益或心理援助热线；必要时伴侣共同咨询。
  \item 一句给男性的话：\textbf{"扛住"不是勇敢，"说出口"才是；你不必为别人的暴力负责。}
\end{itemize}

\section{一页要点}
男性也是性侵受害者，路径与 A 册一致：安全→就医（泌尿外科/急诊）→STI/损伤/伴侣妊娠评估→心理支持。
男性专科要点：生殖道与泌尿道损伤、会阴/盆底疼痛、致伴侣妊娠风险。
开不了口时，一句话"我可能遭了性侵，想检查"就够；找同性别医护陪同、按自己节奏来。




\clearpage
\part*{附篇：旧书整理（1990s 家庭视角）}
\addcontentsline{toc}{part}{附篇：旧书整理（1990s 家庭视角）}

\begin{center}\zihao{5}\itshape
以下 4 章整理自 1990s 旧书，观念内核（沟通、尊重、体谅）仍有效；\
具体数据/手法请以前文现代章为准。
\end{center}

\chapter{旧书整理章·要点速查}

本附篇收录 4 章 1990s 旧书内容，逐章给要点与过时提示：

\begin{itemize}
  \item \textbf{男性生理认知：阴茎·勃起·梦遗}：阴茎大小差异/勃起机制/梦遗的正常认知；\textbf{口语化}，核心（这些都是正常生理）仍成立。
  \item \textbf{男性在性交中的意外}：性交中可能出现的意外与处理；\textbf{实用}，安全提示仍适用。
  \item \textbf{妻子如何面对丈夫阳痿与自慰}：伴侣视角的沟通与体谅；\textbf{夫妻向}，核心（沟通、不指责）仍有效。
  \item \textbf{男性不育}：男性不育的病因与检查；\textbf{部分数据过时}，以现代男科/生殖医学科为准。
\end{itemize}

\chapter{男性生理认知：阴茎·勃起·梦遗（旧书整理）}
阴茎大小有区别


阴茎是男子最显著而且又被谈论得最多的性器官。阴茎疲软状态下的长度，各人之间的差异很大，即使是同一个人，在不同的精神状态和不同的环境下，阴茎的大小也有差别。

一些男子以为自己的阴茎比别人小，从而产生自卑感，甚至担心影响性生活，其实大可不必。因为阴茎的大小对女性的性满足并不十分重要。阴道2/3对阴茎的刺激不太敏感；在性交时，阴茎和阴道的摩擦，刺激了阴蒂、小阴唇、阴道口及阴道的外1/3段，才使女性获得快感。而且事实上由于阴道富有弹性，男性阴茎的大小并不影响女性性交时的快感。因此说，阴茎短小，不用担心，只要阴茎能勃起进入阴道性交，它就是符合标准的。





阴茎的勃起


阴茎正常勃起，甚至在没有性状态下勃起都是正常的，而且是健康良好的征兆。有些男性大约在8岁左右开始增加自然勃起的次数，这种不自主的勃起常发生在没有性念头或性刺激的情况下，可能会在生殖器摩擦到衣服时，感到寒冷时或淋浴时有勃起现象，这些都是正常的现象，企图控制自然勃起可能对以后的性功能造成影响。事实上，这种自然勃起几分钟内就会消失，但你若感到焦虑则可能会使勃起更持久。





夜间射精


夜间射精（“梦遗”）也是完全正常的现象，而且完全无法由意志控制。一般来讲，正常情况下出现的“梦遗”对身体是无害的。因为睾丸在不断地制造精子，前列腺和精囊等会分泌精浆，两者形成了精液，当精液积累到一定程度，就会以梦遗的方式捧出体外，通常青春期的男性，一般十几天可发生一次梦遗，但如果每隔两三天就发生一次梦遗，那就有问题，应及时去医院检查治疗。

\chapter{男性在性交中的意外（旧书整理）}

男性在性交中的意外


男性在性交中的意外死于比女性稍多见。如男子过度兴奋，神经中枢失去抑制作用，最后精脱死亡。又如不顾奔波劳累，纵欲过度，也可虚脱死亡。危机时可用指甲或发卡刺腰部等紧急措施。男子的神经中枢或许立即发生反应，制止令人精力衰竭的痉挛性发作。

也有的人是在连续激烈的运动、工作后，身体疲惫时强行同房。精脱而死，古人云“百里行房者死”。这是因为体内激素大量耗竭，疲劳体虚，进而阴虚阳亢所致。遇到这种情况，可积极救治，如掐人中，轻咬舌头、嘴唇等。另须注意极度疲劳者不可不慎重行事。饱食之后即刻发生性行为也很危险，轻者导致终生消化不良，长期腹泻，重者有生命之忧。

性行为后大汗淋漓，过分贪凉；暑日恣情纵欲，均易引起意外。

（1）激烈的性交可使男子发生尿道炎。

（2）长时间的激烈的性交可能引起尿道创伤，引发疼痛。已有尿道炎、尿道结石、尿道憩室、尿道狭窄等异常均可使尿道损伤。

（3）使用器具、咬伤等可使阴茎皮肤出现创伤，阴茎出现水肿。

（4）包皮过长过紧的男子包皮可无法回到原来位置。

（5）睾丸被挤到耻骨或腹股沟造成睾丸损伤。

（6）阴茎撕裂多发生在阴茎的中段和远端。在损伤时发生“喳”的一声响，然后是剧痛和迅速肿胀、变色和变形，当尿道口流血、尿血或排尿困难时应怀疑有尿道损伤。勃起阴茎受到创伤时可发生阴茎折断；性交急剧改变性交体位时一般会发生阴茎折断。

\chapter{妻子如何面对丈夫阳痿与自慰}


什么是“阳痿”


广义上来说

凡性欲低下，阴茎不能勃起，射精快等性功能障碍统称为阳痿。

狭义的含义

即指阴茎不能勃起进行性交或阴茎虽然能勃起，但不能维持足够的硬度以完成性交，时间已达1年以上。

阳痿是一种复杂而整体效应特别明显的性障碍性疾病。正常性功能也存在着生理波动。在正常性刺激下，反复多次的出现性交失败，方能认为是阳痿。阳痿发病率约为成年男性的10％，其发生常随年龄的增长而增加。





阳痿常见人群


阳痿可以发生在任何男性身上，不分年龄、职业或种族界限。无论是运动员、商人或是工人，都同样有机会患有阳痿。

大多数男性偶尔碰上阴茎无法勃起或难以持续勃起的问题。一般起因是由于过度疲劳、精神过于紧张或饮酒过多。

暂时失去勃起能力，根本不足为虑。然而，倘若这种情况持续长时期或影响房事，便应到医院诊治了。





阳痿的成因


以往，医学界一般都认为阳痿是由于精神或情绪问题所引致，因此以往医生通常只会规劝患者靠自己努力克服这个问题，直至最近才改变了看法。

时至今日。专家相信生理状况引致阳痿占百分之七十五，而由于心理问题引起的则占其余百分之二十五。但也有很多个案是同时由生理及心理因素所引致的。

生理因素

血管发生毛病是最常见引致阳痿的病因。这是指血液流至阴茎的数量不足够（动脉充血不足）或过度流出以致无法持续勃起（静脉血液流量过多）。

动脉充血不足可能是由于多种情况所引致，包括：动脉粥样硬化、吸烟、过度紧张、糖尿病、会阴或骨盆创伤以至血管受损，或骨盆放射治疗。此外，尿道下部受创伤引致阳痿亦占很高的百分比。就五十岁以上的男士而言，动脉充血不足是导致阳痿的主要生理病因。

静脉血液流量过多则是由于静脉机能衰竭所引起，起因是由于平滑肌不够松弛、阴茎神经受损、或阴茎内海绵状横隔膜之结构变更。

因药物而引起的阳痿，通常是由于不同种类的处方药物产生副作用所致。目前有许多药物用以治疗高血压、精神沮丧或精神病，都有可能会扰乱肌肉的惯常行为，因而引致阳痿。

根据美国一项名为Tomhs的治疗轻微高血压研究实验发现。钙通道阻滞剂、p阻滞剂、a-1阻滞剂、ACE抑制剂及利尿药等多种用于治疗神经紧张的药物，都会对勃起机能产生负面作用。

结构性原因是指与生俱来的畸形生殖器官，或是该器官曾受到严重创伤，但这些情况较为罕见。

基于神经性原因引起的阳痿则关乎神经系统机能不良。包括：由于糖尿病或酗酒引致末梢神经产生毛病，脊髓创伤；多发性硬化，和手术引致骨盆部位受损。

基于内分泌原因引致的阳痿并不常见，这包括：

由于性腺机能减退引致睾丸激素下降至低水平；大脑垂体生肿瘤引致催乳激素处于过高水平；和甲状腺素过度及甲状腺素不足。

吸烟及酗酒亦会引致阳痿。香烟引致血管收缩，并增加患上血管疾病和高血压的风险，此两种疾病均会引致阳痿。美国研究报告指出，男性大量吸烟（每天吸烟多于四十支），在梦中竖阳次数最少，而且较难持续勃起。

酗酒会引起末梢神经产生毛病、睾丸机能衰退、影响脑垂体阻碍肝脏功能。因而使血液中的雌激素上升，导致阳痿。

心理因素

正如脑海中经常想念性事会令男性勃起，反之作负面的思想也会妨碍勃起。过分渴望有好表现、负面的自我形象、婚姻不和谐、沮丧、滥交及罪恶感，都会引致阳痿。

男性如患上基于心理因素引致的阳痿症，通常会在梦中或在清晨起床时竖阳，但是面对性伴侣时则可能有心无力。





男性如何受阳痿影响


性功能健全能界定男性在社会上所扮演的角色，而确定其身份。所以，丧失勃起能力会对男士造成极大的影响。

许多阳痿患者在情绪上会感到困扰，而经常表现沉默。他们可能告诉自己：“如果我不能与伴侣有正常的性生活，我便是一个失败的男人及情人。”他们又可能担心：“如果我不能使她满足，她会离开我吗？”这些忧虑可能导致焦虑、情绪低落、悲伤、尴尬、愤怒、羞愧、挫折等等各种感受。





伴侣如何受阳痿影响


由于阳痿造成情感上的痛苦，所以男士常会制造借口，避免与伴侣有性接触，以此忘记病况。不论有意或无意，这样的做法会使女性感到彼此关系冷淡，因而经常对人有所抗拒，并感到寂寞和情绪低落。

有些女性可能害怕自己对伴侣已失去吸引力，或以为受到这样的对待全因为自己曾做了什么。另一些可能担心：“我们的关系是否出了什么问题？”或“他另结新欢么？”因此，男性不与伴侣讨论自己的情况，可能会令伴侣感到忧虑和情绪低落，亦可能令她大发雷霆或感到不安。

可是，通常男性宁愿自行设法应付阳痿的问题，也不愿伴侣施以援手：有些甚至隐瞒伴侣，自行求医或寻求辅导。大多数医生认为。一双伴侣应合作找出对彼此有利的方法，解决阳痿的问题，收效才会大。





阳痿的治疗方法


阳痿是可以安全有效地治愈的。现代科学昌明，男性无须继续沉默忍受痛苦，任由不举破坏性关系。任何人如果怀疑自己患上阳痿，便应延医诊治。泌尿外科医生便是诊治阳痿的专家。

大部分阳痿患者都不愿意求诊，这是由于很多人都不喜欢讨论性生活问题。而且很多人仍然抱着守旧思想，以为阳痿是因心理问题或以为阳痿是不治之症，因此到目前为止。阳痿仍然未能得到正确而全面的诊断和治疗。虽然如此，最重要的是要知道，阳痿是可以治愈的。治疗的方法包括下列各种：

注射治疗：注射药物在患者的阴茎两侧以增加血液流量。

真空器具：乃一个圆柱体的外用气泵，用来促进血液流通至阴茎，刺激阴茎勃起。待至血液流量充足时，随即使用筋套在阴茎的根部，以持续勃起。待性事完毕后，再除去橡皮筋。

人工阴茎：须进行外科手术在阴茎两侧放置一种人工阳具，以便随心所欲勃起。放置的人工阳具有两种形式：一种是可充气的装置。另一种则是半硬的小棒。除非尝试其他方法无效，否则不推荐使用这种方法。

外科手术治疗：通常只适合一些因阴茎或骨盆部位受伤以至血液流通闭塞的患者。这种手术疏通闭塞的范围，令血液流通回复正常。

性心理辅导：可以减轻因阳痿而引起的过度焦虑情绪，从而减少病况的影响和持续时间。这种治疗法多结合其他由医生指导的治疗方法一同采用。

荷尔蒙治疗：可以调节睾丸激素、催乳激素及雌激素等各种性荷尔蒙分泌，对于因为荷尔蒙失调引致阳痿的男性最有效。

口服治疗：可以对少数阳痿男性有所帮助。印第安人及非洲人使用一种从常绿树树皮中提炼的药物Yohimbine作为催情药，已经有一段悠久的历史。这种药物相信有助于加强阴茎的神经敏感，但迄今还没有对此疗法进行大型的研究。

改变生活习惯：例如戒烟便是减轻甚至消除阳痿影响的重要步骤。改变处方药物也有可能解决这个问题。

无论是什么原因，男士必须承认一点，就是阳痿导致夫妻间关系疏离。而在大多数的情况下，只要简单的治疗便可“重振雄风”。所以，若你怀疑自己患有阳痿的话，请勿受传统保守思想的阻碍，因羞于启齿而延误求医！记着：阳痿早日痊愈。夫妻间便可早日完全享受亲密性接触带来的满足和欢乐。到时，你的生活便会更圆满。

你最近可有发觉自己勃起时硬度不足，以致未能完成房事？





自慰不仅仅是性爱的代替物


专家进行的调查发现，自慰不仅是性爱的代替物。1970年，美国专家所作的广泛性行为调查显示，百分之七十二的丈夫，百分之六十八的妻子有自慰行为；其频率，丈夫一年平均自慰二十四次，妻子则只有十次。大多数结婚的男性仍然自慰已是一个事实，问题是，他们为何要自慰呢？

三十五岁的工程师贾某，结婚已八年，对自慰有以下的看法：“有时，性爱变得有压力。它不只是一种身体活动，婚姻已将它变成负担。性爱变成了你对妻子表现爱意的行为，妻子将它看成是有意义的表现。还有，如何使她达到高潮，也增加了压力，在进行性爱时，我差不多有一半时间，是在担心她能否达到高潮，而没有期望本身可享受到乐趣：总想着她的乐趣，妨碍了我得到的乐趣。自慰是最简单的解决方法，不必考虑别人。”问他，他的妻子是否知道他自慰。他回答：“当然不知道。她会认为这是对她的‘背叛’。”那么，他在何时何地做这种事呢？“通常，我在浴室作。这是你唯一可锁上门逗留的地方。那里有很多的抹布，可清除所有的痕迹。”

四十岁的许某，是政府官员，结婚已十三年。他说，自慰是他的必要活动。“一句话，我需要的性爱比妻子为多，我尝试用各种方法刺激她兴奋。她仍然不想一周造爱超过两次。我不想自吹，但即使我在晚上造爱两次后，到了半夜，我又会想造爱。我感到，我的性欲仍然像青春期般旺盛。坦白说，如果我不自慰，我会走上婚外情的道路。”

问许某，他让妻子知道，是否他不想偷偷摸摸进行。他说。“我问她。她是否想看，我希望她看了后产生兴趣，一起进行。她根本不考虑这件事，她说，看电视更有意思。”





自慰是追求新奇的一种反应


专家认为，已婚男性自慰，是“十分正常的”，因其可弥补长期关系中失去的东西，它是对新奇需要的一种反应，并能为性幻想提供了良好的机会。已婚的男人嫖妓，是探索婚姻失去的两种同样的东西——新奇和性幻想。至于自慰是否为伴侣接受，完全视乎她是否担心而定。

一些男性承认即使自慰，他们亦同样追求新奇。“预先知道右手作的事，令人厌烦。”理查坦白说，他改用左手作。体会一下有何不同。五十一岁的经理詹利迈说，“我知道，自慰使我产生不忠的感觉。为何我还作呢？大多数时候，我只和性感的妻子造爱。结婚二十一年来，我们在性生活方面很满足。然而，我偶然也会自慰。我作这件事，完全是为了改变花样；我感到它的确不同，它结束得快些，能吸引更大的注意，显得粗犷，和造爱的感觉不一样。毫无疑问，我喜欢实际的性交。不过，偶然自慰，性的改变能带来一些刺激”。





丈夫自慰通常害怕妻子当场发现


虽然男性自慰，却怕妻子当场发现。二十七岁的中学教师杜格拉斯说，“当时，我真想死。一小时前，我还和她性交过。我只是想看看裸女照片，并对着相片自慰起来；没想到，她到洗手间拿纸巾，一下撞见了我。”

为何男人在造爱后还想对着裸女照片自慰呢？这只有一个解释：渴望。和女性看时装杂志一样，纵然她们身上穿着漂亮的衣服，依然要去看那些昂贵的衣服，纯粹出于渴望而看的。

杜格拉斯解释他所做的事，“我感到自慰是性幻想的最佳时候，选看性感的女人照片，有助我进行幻想。虽然我在一小时前性交过，但是，我接着又有性幻想的渴望，所以我作了那样的事。”

杜格拉斯对妻子怎样说，她的反应怎样呢？“我气急败坏地向她道歉，但是，这是一场噩梦。她流出眼泪，猝然离开，称我是性反常者。她似乎认为，这是精神上的不忠，和肉体上的不忠一样坏。至于她担心的是，这显示，她不能够满足我：老实说，我不能相信她竟然会大惊小怪。很多男人都做这种事，我不关心他们是否告诉他们的妻子。这是你可完全控制的性形式。”

杜格拉斯承认，自从那天后，他的婚姻和以前不一样了。“我感到，她似乎在提防我找出裸女照片，开始对着自慰，”他说，“我提议性爱，她对我说些讽刺的话，例如：‘你的阴茎有空了吗？’”他只好暂时停止性爱。“有一夜，我们互相爱抚一阵后，我诱导她说出了结婚后，她也自慰过。她说我到外地时，也自慰过两次。现在她承认了，我们都一样。”

一些妻子说，知道丈夫自慰；对这样的事无所谓。“尽管你很爱你的伴侣。你们的性欲不总是一致的，利用自慰来解决问题，未必不是一件好事。”一位女性说。





快速的性松弛


专家指出，婚姻中的自慰，完全是正常的现象。在婚姻中，实际上有两种性本能发生作用：一种是两个人的性爱，这是全面的肉体性爱，需要时间、情绪，和花费很大的精力；另一种是快速的性松弛，不到一分钟的时间便解决，这便是自慰。当你有性的冲动，而对方没有性的兴趣，你偶然地用自慰来松弛，比服用安眠药更能有效助你入睡。

\chapter{男性不育（旧书整理）}
男性生殖系统肿瘤


男性生殖系统的肿瘤，能影响生育的有阴茎癌、尿道癌、睾丸肿瘤、精囊囊肿、附睾肿瘤以及前列腺增生和癌。

阴茎癌多先在阴茎头及冠状沟发病。在病因方面，除考虑包茎和个人卫生习惯较差以外。可与不洁性交，尤其是梅毒有关。阴茎癌无论是处于早期还是晚期，实际上患者从发病之初起就因勃起时疼痛或出血而中止性交。本病患者早期常不到医院诊治，所以一般到就诊时已大多为晚期，此时大多数患者只能被迫进行治疗性阴茎切除。如果只作局部淋巴结手术和化学疗法，虽然是治疗上的改进，能保留阴茎，但尿道狭窄、难以愈合的坏死和疼痛，是化学疗法普遍的合并症。佩罗尼氏病是阴茎海绵体上发生纤维斑块硬结，尚不能十分肯定是肿瘤。当阴茎勃起时常向一侧弯曲，严重者不能性交，时常合并阳痿，影响到生育。在治疗上有时不得不切除白膜，植入假体。

睾丸肿瘤是临床常见病，其中非生殖细胞瘤很少见，94％来源于生殖细胞，其中以精原细胞瘤与畸胎癌为常见。无论是哪一类型的睾丸肿瘤。病情轻者，睾丸组织会受到损害，严重时能危及生命。但双侧睾丸同时生长肿瘤者极少，如果是一侧睾丸肿瘤，暂时是不会影响生育的。但在睾丸肿瘤治疗过程中，有时因做了腹膜后淋巴结清除手术，会引起射精障碍而不育。应用化学疗法，有时能引起睾丸萎缩和脱发，并且需要经过很长时间，睾丸萎缩和精子缺乏才能逐渐康复。

尿道肿瘤尤其是癌，临床不多见。无论肿瘤生长于尿道任何部位，发病不久即会因充血而发生疼痛、出血、排尿痛与性交痛，均伴有尿道狭窄。其所以影响生育主要在于性交困难或不能性交。精囊肿瘤主要是囊肿，极少见且多为单侧发病，生育并非完全不可能。精索肿瘤，几乎全为胚胎性，大多为良性。如果闭精与输精管道正常，对生育不会有太大的影响。附睾肿瘤虽然十分少见，但由于附睾肿瘤的存在会影响精子质量、数量和输送。在一定程度上会影响生育。双侧附睾同时发生肿瘤者甚少。





男性生殖器官炎症


男性生殖器官的炎症与泌尿系统器官的炎症有着极为密切的关系，很多情况是互为因果的。前列腺炎尤其是慢性前列腺炎，是成年男性常见的疾病，该病在病程的某一阶段，会出现阳痿、早泄和神经衰弱等症状，无疑会影响生育。所以如发现有阳痿、早泄等症状时，首先要确定有无前列腺炎症。无论是细菌还是病毒所引起精囊炎和前列腺炎，都可抑制生精功能。尤其是精液中含有细菌，即或无任何临床症状，也能引起精细胞分解、粒细胞中毒和精子凝集等现象。

睾丸炎常为阴囊炎的并发症。睾丸炎常见的原因是附睾炎、腮腺炎和梅毒。实际上，无论原发于睾丸的炎症或来自其他系统的炎症，均损害睾丸。流行性腮腺炎引起的睾丸炎，常发生在腮腺炎一周之后，大多为单侧。如为双侧可因睾丸萎缩导致不育，此时可见生精上皮的严重损害。单侧睾丸受损害后仍然有生殖能力，若双侧受损害，则能造成少精症或无精症。由淋病菌感染的睾丸炎与附睾炎，能导致睾丸功能丧失和附睾部分梗阻或完全阻塞。

附睾炎，尤其慢性附睾炎，常与慢性前列腺炎和慢性精囊炎同时存在或先后发生。附睾炎常因结核杆菌和淋病双球菌感染而发病，但非特异性感染并不少见。在炎症性病变发展过程中，附睾管腔可因瘢痕形成变细变窄，有可能发展成管腔完全闭塞。如果病变为双侧，无疑会造成不育症。包皮炎与龟头炎属常见病。病期因性交困难而不育，一般病程较短，炎症愈后对生育无大影响？但应警惕淋病与梅毒在阴茎头部所出现的症状。

一些无症状的生殖系感染也可造成不成熟精子的早期脱落，如生殖道内的支原菌属感染。





精子输送障碍


精液不能射入配偶生殖道内是指男性的精子输送通道正常而精液不能射入女性阴道内。通常有男性和女性两方面的原因。男性方面的有以下原因，可对症治疗。

1．先天性畸形，常见的畸形是一侧或双侧输精管缺如或闭锁，也有半附睾缺失者。

2．炎症，男性生殖系统的急、慢性炎症均可引起精子输送通道管腔狭窄或阻塞。常见的炎症有附睾、输精管、精阜和前列腺等急、慢性炎症。在慢性炎症中，要特别注意结核杆菌引起的特异性炎症。此外，还应警惕由性病引起的生殖系统炎症。

3．肿瘤，大多为良性肿瘤，如附睾腺瘤样肿瘤、平滑肌瘤和囊肿、精索腺瘤样肿瘤、前列腺癌、阴茎癌等。上述肿瘤有的来自精子输送通道自身，有的来自精子输送通道周围器官，都可造成通道管腔内堵外压，影响精子通过。

4．外伤，阴囊或会阴部外伤，可损伤附睾和输精管，引起管腔部分或全部断裂，愈合的瘢痕也会影响精子通过。

5．其他病症，如巨大的腹股沟病、睾丸扭转等也可引起精子输送通道梗阻。





遗传性疾病


患以下遗传性疾病可发生不育。

1．先天性睾丸发育不全，此病较为常见，是一种遗传性疾病，是因染色体异常伴有睾丸发育不良导致男性不育的典型病征。

2．XYY综合症，此病在男子中占0.2％，染色体核型为47XY或XYYY。临床表现仍为正常男性外观，只是身材高大，可为无精子或少精。性激素水平基本正常，或升高。

3．特纳氏综合症，本病患者大多数的染色体核型仍是46XY；仅少数是染色体嵌合。临床表现为身材矮小，耳朵位置低，盾状胸，肘外翻，有心血管系统的畸形。性机能延迟发育。隐睾及造精功能障碍，血清中雄激素水平低。多数丧失生育能力。

4．性逆转综合症。本病的临床表现与先天睾丸发育不全综合症相似，外表为男性，但睾丸小，无阴毛或稀少。染色体核型为46XY。

5．单纯支持细胞综合症，发生率约为3％。其特征是外观上无先天畸形表现，但睾丸中缺少生精细胞，无精子，血清促性腺激素水平升高，但睾酮水平正常。

6．纤毛呆滞综合症，此病家族史为近亲结婚多，发病率高。临床特点是内脏反位。支气管扩张和鼻窦炎。此病出现纤毛运动障碍，所以容易并发感染。纤毛运动障碍也可以发生在精子，引起精子尾部纤毛运动失常导致男性不育。

7．范康尼氏综合症，此病往往在一个家族内发病，临床表现是出生时身材矮小，骨骼发育畸形，如拇指缺损，指趾畸形及色素沉着等；生殖器发育不良，性功能不全。血液方面因骨髓中血细胞形成不良、血小板形成不良而表现为全血细胞减少。

8．Prader Will氏综合症。此病从乳儿期开始肌力低下、肥胖、身材矮小、肢端短小、智力低下、性腺发育不良、睾丸不下降、小阴茎等。

9．多腺性内分泌不全症，小儿和青年发病多，可见家族性发病，属于常染色体隐性遗传性疾病。临床表现为恶性贫血、重症肌无力、肺纤维症、慢性活动性肝炎、性腺功能低下。最近发现患者血中有对产生类固醇激素细胞作用的抗体。此种抗体与性腺功能低下的临床表现有关。

10．劳穆比氏综合症症，本病可能与妊娠时接受放射性物质或化学物质致使间脑、下丘脑障碍有关。临床表现为肥胖、视网膜色素变性、智能低下、多指（趾）畸形、尿道下裂、生殖器官发育不良等。

11．勒里斯氏综合症，由于髂总动脉狭窄甚至闭塞，致使下肢和男性生殖器官供血不足。早期症状表现为阴茎不能勃起，两下肢严重疲劳及肌肉萎缩。晚期出现性功能完全丧失。





自身免疫性不育


男性不育的病因很多。对抗精子的免疫反应是其中之一。精子是一种自身抗原，可诱发机体产生抗精子抗体。在正常情况下，精子抗原并不与自身免疫系统相接触，因而不会引起自身免疫反应。在正常睾丸中存在有坚固的血睾屏障，可以保护精子抗原免受攻击。由于某些原因，精子抗原与自身免疫系统相接触，便可引起体液免疫反应和调理免疫反应。

抗精子抗体可以存在于男女任何一方或双方，不育夫妇中，由抗精子抗体所致的免疫性继发不育占15％～30％。

人类的抗精子抗体，一种是精子凝集抗体，可以引起精子凝集；另一种是精子制动抗体，可以使精子丧失活动能力。目前仍不十分清楚不育与抗体之间的因果关系。一般认为产生抗精子抗体的原因与下列情况有关：输精管道阻塞、睾丸炎症和损伤、附睾感染、附属性腺的感染等，特别是输精管阻塞。





其他因素不育


1．烟、酒的影响：吸烟对人体健康是有害无益的。吸烟对性功能的影响，也逐渐地被科学界所重视，香烟的烟雾浓缩物以及大量吸烟者的尿提出物中含有诱发细胞畸变和阻碍淋巴细胞合成脱氧核糖核酸的物质，致使吸烟者末梢淋巴细胞染色体的易位。精子的生成过程需要脱氧核糖核酸与蛋白质的合成，因此，吸烟对精子的生成、成熟以及畸形精子的比例都有明显的影响。吸烟者精液中畸形精子比例都远远高于不吸烟者。每天吸烟20支以上者比吸20支以下的精子畸形发生率显著增高。有一部分器质性阳痿的患者是由于动脉血管受损引起的。吸烟乃是引起动脉血管受损，尤其动脉硬化的常见的危险因素之一。

酒精是一种性腺毒素，过量或有长期饮酒嗜好者，可引起性腺中毒，特别是对睾丸，表现是血清睾酮水平降低，出现性欲减退、精子畸形、阳痿等性功能障碍。慢性酒精中毒能严重地损害睾丸间质细胞，造成不能正常地产生雄性激素和精子，并使体内合成睾酮的三种酶的活性受到严重影响。每天平均饮烈性酒250毫升以上，持续2～5年，可以导致严重阳痿。有饮酒嗜好者一旦发现自己有早泄、阳痿、遗精或性欲减退等情况时。必须下决心戒酒。

2．放射性物质：放射性物质都会放射出对人体有害的射线。不仅能降低甚至破坏造血系统的功能，还能损害睾丸的造精功能，且具有明显的敏感性。睾丸的生精细胞在受到放射线照射以后。会出现萎缩以至完全丧失造精功能。小剂量或偶尔接受照射，造精功能也会受一些影响，但是可以恢复的。如果大剂量的照射可能造成不可恢复的生精细胞损害。经常接触放射线的工作又不注意加以防护，也会引起相当严重的睾丸生精细胞的损害或者造成生精细胞中染色体的畸变。为了预防对造精功能的损害，必须严格地执行防护保健措施和有关规定，并应定期进行健康检查。

3．维生素、微量元素的影响：维生素缺乏，特别是缺乏维生素A、B、C、E等可以影响精子发生，引起男性不育症。如维生素A缺乏，可使精子发生能力减弱：维生素B缺乏会影响垂体功能而降低生育能力；维生素E缺乏能造成睾丸损害，但不影响间质细胞；维生素C对防止精液凝集有一定作用。因此，在治疗男性不育症中适当地补充维生素A、B、E和C是必要的。

据研究，缺乏锌可导致不育，用锌制剂治疗某些男性不育症，取得了一定的效果，可能与锌能影响精子代谢和性腺发育有关；微量元素能明显影响精子存活率、活动度和穿透宫颈粘液的能力；在分析铁与精液关系时，发现铁和精液中精子密度有明显关系；提示铁可能与精子发生有关；一般认为镁和锌作用相似，但镁不属于微量元素。

4．甲状腺功能低下：当甲状腺激素缺乏时，全身代谢活动下降，以至引起临床综合症，称为甲状腺功能低下。男性可出现性功能低下、性成熟推迟、副性征落后、性欲减退、阳痿和睾丸萎缩。在女性可有月经不调、经血过多或闭经，一般不孕。无论对男女病人的生育都会产生影响。

5．慢性呼吸道疾病：呼吸系统疾病中如慢性支气管炎可以造成男性不育。杨氏综合症患者有慢性支气管炎或支气管扩张病史，临床表现为无精子或少精，但睾丸体积与血清中促卵泡激素和促黄体激素水平均正常。临床上叫做与呼吸道疾病有关的阻断性无精及少精症。还有卡塔格内氏综合症，这种不育病人的特点是在儿童时期即患慢性呼吸道疾病，其精子尾部纤毛运动不正常，精子向前游动的能力弱。有的只在原处移动或旋转移动。在电镜下可以看到精子尾部结构有异常改变，同时发现在呼吸道粘膜上的纤毛细胞结构也发生异常，所以认为这种不育症是因为纤毛细胞异常所引起的，也叫做纤毛呆滞综合症或者男性鞭毛异常症。

男性不育的治疗主要针对提高精液质量，消除影响精液质量的因素（如感染、免疫、精索静脉曲张等），以及解除精道梗阻等。治疗手段包括药物治疗、手术治疗和辅助生育技术治疗。

药物治疗：内分泌治疗主要适用于性腺功能低下、少精子症、精索静脉曲张术后及特发性不育等。其他不育因素的治疗包括生殖道感染、免疫性精子活动障碍、高催乳素血症、精液不液化、性功能障碍，以及其他全身性疾病。

手术治疗：外科治疗的先决条件是睾丸具有生精功能。外科治疗不育的范围主要为：纠正外生殖器发育反常、勃起功能障碍、改善生精功能以及解除精道梗阻。





中医治疗男子不育


中医典籍中所述及的“五不男”，见于明万全《广嗣纪要·择配篇》，主要是指男子生殖器官先天发育畸形或其他病变影响生育者。即天、漏、键、怯、变五种病症。所谓天，是天宦的简称，即指男子生殖器短小若无，不能生育，指遗精、滑精，即男子精关不固，常自遗泻而影响生育者。键，即男子阴茎被割掉者，以致不能够生育。怯，指阳展不举，即男子由于阳痿而影响生育者。变，指生殖器官两性畸形而影响生育者。

中医古籍《辨证录》曾记载：“凡男子不能生育有六病，六病何谓？”

“一精寒，二气衰，三痰多，四相火盛，五精稀少，六气郁。”

（1）精寒是指下焦虚寒，命门火衰，排出精液温度低，有的形容“冷如冰铁”，难于使女方受孕。

（2）气衰泛指脏腑机能不强，或指体内富有营养的精微物质不足，此处尤其是指肾气不足，肾气衰则肾精产生的内在动力不足，影响生育。

（3）痰多与脾、肺二脏有关。中医认为：“脾为生疾之源，肺为贮痰之器”；“脾为气血生化之源”；若痰湿蕴郁脾胃，必定导致真气不足，精气亏耗，同样影响生育。

（4）“相火盛”指肾阴亏损，虚火亢盛，又称命门火旺。由于阴虚火旺，出现精泄的病变。肾为阴脏，内藏水火必须保持相对平衡。若肾水亏损，则肾火偏亢，出现性欲太过、遗精、早泻等。

（5）精稀少，在中国古代医籍《清病源候论》中称为虚劳精少，指性交时泄精少，甚至只一二滴，影响生育。由于先天不足，或房事不节，劳心过度，以致耗损精气。

（6）气郁，郁症之一，由于情志郁结，肝气不舒所致，气郁可导致血瘀，造成阳痿、不射精等症而致不育。





治疗男性不育的常用中成药


龟龄集

有补肾壮阳的功效，适用于肾气虚弱、阳事不举、婚久不育、阴寒腹痛、腰膝疲软无力等症。散剂，每服1.5克，每日1～2次，温开水送服。

大造丸

有补肾滋阴、泻热潜阳的功效，适用于肾亏损、阴虚阳亢、骨蒸劳热。形体消瘦，不育症等，蜜丸9克，每日二次。

通关丸

有滋阴降火，引火归元的功效，适用于肾阴亏耗，相火妄动，滑精白浊，小腹酸痛、缩阳症，形容憔仲，头晕耳鸣等症。水泛为丸。每服9克。每日二次，空腹温开水送下。

参茸卫生丸

有培补肝肾，益气养血、健脾益肾的功效，适用于身体衰弱、精神萎靡、腰膝酸痛、四肢无力、健忘失眠、梦遗滑精、不育、妇女崩漏不止，赤白带下等症。蜜丸每服一丸，每日二次。

海马保肾丸

有滋阴益气，补肾助阳的功效，适用于肾气虚寒引起的阳事痿弱，精神萎靡、四肢无力、梦遗滑精、潮热劳嗽、盗汗失眠、腰膝疲软、目暗耳鸣、嗜卧少食、婚久不育等症。蜜丸每服6克，每日二次。

金匮肾气丸（桂附地黄丸）

有补益肾阳的功效。适用于肾阳不足，腰膝冷痛、阳事不举、小便不利或多尿等症。蜜丸每服一丸，每日二次。

知柏地黄丸

有滋阴降火的功效，适用于肾阴亏虚，相火妄动引起的骨蒸劳热、虚烦盗汗、腰背酸痛、梦中遗精等症。蜜丸每服一丸，每日二次，温开水或淡盐水送下。

鱼鳔丸

有滋阴补肾，添精益髓的功效。适用于气血虚弱、肾水不足引起的腰膝疲倦、梦遗滑精、阳痿早泄、精亏不育、失眠健忘等症。蜜丸每服一丸，每日二次。

锁阳固精丸

有补肾壮阳、温肾固精的功效。适用于肾虚引起的阴囊潮湿、梦遗滑精、目眩耳聋、腰膝疾痛、四肢无力等症。蜜丸每服一丸，每日二次，以淡盐水或温开水送下。

五子衍宗丸

有滋补肾水，添精补髓功效，适用于由阴虚气亏，肾水不足引起的阳控、遗泄及不育。蜜丸每服一丸，日服二次。

六味地黄丸

有滋补肝肾的功效，对肾阴不足引起的不育症，并伴有身体消瘦、头晕目眩、遗精、滑泄、舌燥咽干等症有效。蜜丸每服一丸，日服二次。

人参鹿茸丸

有滋肾益气、补血生精的功效，适用于由肾精亏损引起的气血两亏，精神不振、目暗耳聋、遗精、腰腿酸软、男子不育等症。日服1～2次，以温1升水或黄酒送服。

三肾丸

有滋阴益气、补肾壮阳的功效。适用于因肾气虚弱、阴虚火衰引起的阳痿不举、腰腿酸痛、精神疲倦、食欲不佳、婚久不育等症。蜜丸每服一丸，每日服1～2次。

大补阴丸

有补肾水、降虚火的功效，适用于肾阴不足、相火偏旺引起的耳鸣耳聋及阴虚发热、盗汗、造精不育等症。蜜丸每服一丸，日服1--2次。

左归丸

有滋补肾水、调养精血的功效，适用于肾精亏损引起的腰腿疲软、自汗盗汗、遗精、神疲口燥、头目眩晕、眼花耳鸣、不育等症。小丸每服10克，大丸每服二丸，每日二次。

右归丸

有温补肾阳的功效，适用于先天不足、命门火衰引起的脾肾虚寒、劳伤过度，出现腰酸腹痛、阳痿不举、畏寒怕冷、小便自遗、婚久不育等症。每服一丸，每日二次，用淡盐水送下。

全席丸

有温肾壮阳、固精益气的功效，适用于肾阳亏损引起的精神衰惫、神志不宁、头眩耳聋、遗精盗汗、面色萎黄，婚久不育、腰膝无力及妇女血亏，崩漏带下等症每服一丸，每日二次。





不育夫妻的按摩及功法治疗


方法一：

（1）仰卧位：在腹部用摩腹法操作约10分钟，重点在下腹部，并按揉关元穴，性选用气海穴，女方则选子宫穴，治疗应使腹部深处发热为宜。

（2）仰卧位：以手掌按揉大腿内侧，反复3遍，力量稍重。再选穴位按揉，女性用三阴交穴，男性用太溪穴，稍用力按揉至酸胀并保持约2分钟，然后擦大腿内侧，以热透人深部为宜。

（3）俯卧位：沿整个背腰脊柱内侧，用手掌揉推数遍，力量稍重。再分推腰部10遍，然后按揉腰部肾俞穴，以点按出酸胀的感觉后，缓缓施力，保持按揉3～5分钟。

（4）俯卧位：用手掌分别横擦腰部肾区，斜擦骶部八骨穴，均以热透入身体深部为佳。再施用拿风池颈项法和拿肩井法，共约2分钟。

男性不育患者可于早晨起床前和晚上入睡前各行一次指摩会阴急脉法。其方法为侧卧，舌舐上腭，意守丹田，双腿屈曲，用右手中指指端逆时针方向旋转按摩会阴50～100次。继而仰卧，再以双手中指指端同时逆时针方向旋转按摩急脉50～100次。急脉在腹股沟（俗称大腿根）中点，正置阴毛之边际处。

方法二：

男性不育患者还可以练金刚坐、闭气按摩法。

（1）金刚坐法：盘腿而坐，右足跟抵会阴穴，左足跟抵右足解溪穴，涌泉贴靠在右小腿附阴穴，顺项提顶，竖起脊梁，腰不可软弱。身不可倚靠。两手置扶于两腿血海穴，闭目垂帘，去除心中杂念，待气定神敛后，行闭气按摩法。

（2）闭气按摩法：以鼻吸气闭之，用两掌相搓擦极热，急分两手上下搓摩后腰两边软处，右足跟随手下搓与会阴穴发生碰撞，一面徐徐放气从鼻出，周而复始24遍为1度。可行1～3度。然后变金刚坐法。用上述操练法，使肾脏与生殖系统的废物迅速得到排除和贯通，同时强化呼吸，故亦称“贯通补气法”。

（3）收功法：依次觉察从头至足放松3遍，由“功能态”恢复到正常生理状态。继而缓缓摇头晃腰，左右、前后的摆动脊柱、慢慢舒臂伸脚，搓搓手脸，下坐收功。

注意事项：

1.练功初期有流汗流出；脊背局部发胀的感觉。随功夫的深入，有“腰肾如汤煎”，会阴穴有温热气之感。

2.练功要循序渐进，时间要掌握在20分钟左右，室内应自感不冷，搓摩时要注意命门穴。

3.练功期间要节制性生活，可促进练功效果提高并保证精子成熟。

方法三：

推拿对男性不育是一种很好辅助治疗手法。推拿力量要因人而异，用力均匀，争取先轻后重再轻三个步骤。

（1）束腹法：仰卧位。以两手四指对置在患者脊柱正中线两旁，两手将腹两侧抱拢，由后向腹前摩动，经后腰、侧腹至腹正中线止，反复作束腹式弧形摩动3～5分钟。

束腹时应使两手掌心与皮肤表面贴紧；束腹至腹正中线，两手将腹正中的皮肤肌肉捏紧，并向上拿提；操作时有腰肌向前牵拉感，有腰部轻松及增益精神的感觉；有腹胀者慎用。

（2）下腹横摩法：平卧位。以手指或两手四指掌侧并置于患者下腹部左或右侧，即髋骨内上缘前方约平脐四指处，向腹正中平行摩动，反复横摩5～10分钟。

操作时手掌心微向上抬高。手指掌侧平放摩动；横摩时用力宜稍重；便秘患者左侧横摩用力加重。

（3）命门穴按摩：俯卧位。以食指背屈或拇指掌侧揉腰部之命门穴2.5分钟。再以命门穴为中心，一手掌心置其上，作旋转按摩1～2分钟。

指揉时应轻揉有力，揉动时嘱患者腰背肌放松；操作时防止擦伤皮肤。

方法四：

精子畸形症的按摩：提睾功法和阴囊搓擦法，这套功法不但对治疗精子畸形症有效果，对其他精子、精液异常、性功能障碍患者也都有作用。

（1）提睾功法：作用：能提高脏腑的生理功能，起到固肾益气、壮阳补精作用。功法：站立，两足分开与肩同宽，膝盖微曲，站立宁神。两手由下往上抄，掌心朝上，两手随势上托提气。意念由涌泉穴拔气上提，经两足内侧至阴廉，复会阴，将会阴元气上提人腹，至气海、下脘即可。自涌泉至下脘，吸1/2的气（不要吸足）接着慢慢吐气，同时将舌放平，两手掌翻转，中指相对，掌心皆朝下。随着呼气下推。只吐1/4，足趾要抓地，以逆式呼吸法。吸气时收腹，提会阴。用意念将气由涌泉，经足内踝直提至会阴入腹，由下丹田上至中脘，入膻中。两手掌亦随之提到膻中。身体要向前倾30度。练时要通过意念来带动气功的运动。初练时，只可引气至中脘或鸠尾。待功夫练熟了，可提至膻中或天突，甚至可拉到百会穴。初练时，切不可勉强，以防冲气出弊。继之，再反掌。渐渐下推，气由膻中逐节下送至涌泉穴。意念之气由涌泉穴人地3尺，每人1尺，吐一口气。吐气时，足趾放松。

（2）阴囊搓擦法：阴囊搓擦是“提睾功法”的辅助动作，但相当重要。其方法是以手心搓擦阴囊，一手将阴茎上翻，并将两只睾丸推向上面腹股沟内：再以另一手掌心轻轻上下搓揉36次。擦毕，两手交换，以微热为度，每次擦5～10分钟，不要太重，当心擦破囊皮。

方法五：

死精过多的夫妻按摩：

（1）脐旁横摩法。用手掌或食指、中指、无名指指面于脐旁有节律横向抚摩，每分钟120次。

（2）下腹横摩法，用手掌或手指指面于气海、不门、关元穴有节律横向抚摩，每分钟120次。

（3）揉命门法。用大鱼际、手掌根部或手指吸定命门穴，做轻柔的回旋揉抚，每分钟120～160次。

手足按摩：

（1）掐点手部掌侧命门点；按揉背侧全息穴肾点，下腹点。

（2）点揉足部底侧涌泉穴；背侧公孙穴。

（3）重按手部反应区肾区。生殖区及生殖腺区；掌腕部摩推。

（4）点揉足部反射区生殖区及生殖腺区，肾区；踩足底；擦足心。


\end{document}
